\section{An intrinsic quadratic transgression class modulo full first-order exactness}
\label{sec:r11-intrinsic-quadratic-transgression}

\paragraph{Purpose and scope.}
OR makes a quadratic actual-zero presentation full first-order
exact in finite residue rank at most two.  This section records
the resulting intrinsic quotient interface.  It is a formal
consequence of that proved analytic construction and the retained
QP identities; it does not posit a new primitive or approximation
theorem.  The quotient is the full degree-one module modulo
first-order exact terms.  It is not called Aeppli cohomology:
the representatives below need not be \(D_\delta\)-closed.

Keep the original finite \(K/\mathbf Q_p\), its original norm,
and the project's good, Hausdorff, boundaryless quasi-smooth
spaces with their ordinary topology.  All scalar functions below
belong locally to the fixed finite continuous PL/PS closed-cell
class of OR/QP.  The words ``locally'' and ``sheaf'' permit an
ordinary cover of such finite presentations; they introduce no
uncontrolled variable-PS rule or infinite formal sum on one germ.

\subsection{The existing modules and the quadratic image}

Here \(\mathscr B^{a,b}=B^{a,b}_{\rm Mih}\) is exactly the
genuine analytic Mihatsch delta-form sheaf, in total bidegree,
fixed in the opening of the archived QA source.  This is an
existing actual analytic quotient with all actual moment
refinements.  It is not a new sheaf inferred from the sought
arbitrary-background differential calculus.  The precise archived
inputs are QA8--QA9: they define
\(d'B^{0,1}+d''B^{1,0}\) as an ordinary image sheaf,
state that its full derivative \(D_\delta\) is zero, and
explicitly use the already defined map
\(D_\delta:B^{1,1}\to B^{2,2}\).  QA12--QA13 put the
quadratic transgression and its derivative in those same sheaves;
QP4--QP5 extend that fixed scalar class from PL to the present
finite continuous PS class.

In particular \(\mathscr B^{0,1},\mathscr B^{1,0},
\mathscr B^{1,1}\) contain the \(d''u\), constructed
covectors, full Cartier terms and regular energies used by PR.
Their existing full derivatives have nilpotence, anticommutation
and ordinary restriction compatibility.  We set
\(\mathscr C^{2,2}_{\rm act}=\mathscr B^{2,2}\).
This notation neither defines an action of an arbitrary background
on \(B^{2,2}\) nor identifies \(B_{\rm Mih}\) with the
algebraic GK sheaf.  Those separate interfaces are not assumed.

Take images and quotients in sheaves of real vector spaces:
\[
 \begin{split}
 \mathscr J^{1,1}
    &=\operatorname{im}_{\rm sh}
          (d':\mathscr B^{0,1}\to\mathscr B^{1,1})
      +\operatorname{im}_{\rm sh}
          (d'':\mathscr B^{1,0}\to\mathscr B^{1,1}),\\
 \mathscr Q^{1,1}&=\mathscr B^{1,1}/\mathscr J^{1,1}.
 \end{split}
\tag{IQ1}
\]
Thus a section has zero class precisely when it is locally full
first-order exact.  This is a sheaf quotient, not the assertion
that every such section is one globally exact expression.
Nilpotence and anticommutation give
\[
 D_\delta(d'\beta+d''\zeta)=0,
       \qquad
 \overline D:\mathscr Q^{1,1}\longrightarrow
                                  \mathscr C^{2,2}_{\rm act}.
\tag{IQ2}
\]

Let \(\mathscr P_2\) be the sheaf of finite real quadratic
presentations in the stated scalar class, with the usual real
bilinearity and symmetry of \([h][k]\).  Equivalently one may
sheafify the corresponding symmetric-square presentation space.
The coefficients are real constants, not arbitrary scalar functions.
Put
\[
 \begin{split}
 P&=\sum_\nu c_\nu[h_\nu][k_\nu],\\
 \kappa_2(P)&=\sum_\nu c_\nu
                          D_\delta h_\nu\wedge D_\delta k_\nu,
       \qquad
 \mathscr K_2=\operatorname{im}_{\rm sh}\kappa_2.
 \end{split}
\tag{IQ3}
\]
A section of \(\mathscr K_2\) has local finite presentations;
it need not have a single finite presentation on the entire open
set.  QP supplies the real linear transgression expression
\[
 \begin{split}
 \alpha_P&=\tfrac12\sum_\nu c_\nu
             (h_\nu D_\delta k_\nu+k_\nu D_\delta h_\nu),\\
 u_P&=\sum_\nu c_\nu h_\nu k_\nu,\qquad
 A_{0,P}=\tfrac12D_\delta u_P,\\
 Z_P&=\tfrac12\sum_\nu c_\nu
      (d'_Ph_\nu\wedge d''_Pk_\nu+
                       d'_Pk_\nu\wedge d''_Ph_\nu),\\
 \alpha_P&=A_{0,P}-Z_P,\qquad
 D_\delta\alpha_P=\kappa_2(P).
 \end{split}
\tag{IQ4}
\]
The ordinary scalar identities already make IQ4 compatible with
the stated bilinearity and symmetry of the presentation space.

\begin{theorem}[Canonical germ class in residue dimension at most two]
\label{thm:r11-intrinsic-quadratic-germ-class}
Let \(x\) be any original point, in arbitrary ambient dimension,
with \(s_K(x)\le2\).  There is a canonical real linear map
\[
 \tau_{2,x}:\mathscr K_{2,x}\longrightarrow\mathscr Q^{1,1}_x,
       \qquad
 \tau_{2,x}(\kappa_2(P))=[\alpha_P]=-[Z_P],
\tag{IQ5}
\]
independent of the finite original presentation of its input.
It satisfies
\[
                          \overline D\,\tau_{2,x}
                                      =\operatorname{id}.
\tag{IQ6}
\]
Its image is exactly the subspace of quadratic transgression
classes.  On that image \(\overline D\) and \(\tau_{2,x}\)
are inverse isomorphisms.
\end{theorem}

\begin{proof}
Take two finite germ presentations \(P,P'\) of the same actual
quadratic curvature.  Shrink ordinarily so that both are defined,
with one finite union unit tuple \(f\) and one common finite
closed-cell subdivision.  Their difference is an actual zero:
\(\kappa_2(P-P')=0\).  FM's bound gives
\[
                         s_f(x)\le s_K(x)\le2.
\tag{IQ7}
\]
OR applied to this entire difference constructs, on an ordinary
neighborhood, the full identity
\[
                  \alpha_P-\alpha_{P'}=d'\beta+d''\zeta.
\tag{IQ8}
\]
Consequently their classes in IQ1 agree.  This proves
well-definedness for germs, including presentations originally
given in different actual unit coordinates.  Real linearity
follows from the linear expression IQ4 after taking a finite
common refinement.

The full singular term is legitimately exact in this quotient:
\(A_{0,P}=d'(\tfrac12d''u_P)\).  Hence
\([\alpha_P]=-[Z_P]\).  This is an equality of quotient classes
using the complete term \(A_{0,P}\), not a claim that the full
polar has a smooth or continuous double primitive.  Applying
IQ2 and IQ4 gives IQ6.  Every quadratic transgression class is
represented by some \(\alpha_P\), so it lies in the image
of \(\tau_{2,x}\); IQ6 proves the inverse assertion.
\end{proof}

\begin{theorem}[An ordinary sheaf morphism on analytic surfaces]
\label{thm:r11-intrinsic-surface-transgression-sheaf}
On an ordinary open set of pure dimension at most two, the germ
maps above assemble to a canonical morphism of sheaves of real
vector spaces
\[
                   \tau_2:\mathscr K_2\longrightarrow
                                    \mathscr Q^{1,1},
              \qquad \overline D\circ\tau_2=\operatorname{id}.
\tag{IQ9}
\]
It commutes with ordinary restriction and is independent of all
actual coordinate refinements used to represent its input.  More
generally the same assertion holds on any ordinary open set on
which \(s_K(x)\le2\) at every point.
\end{theorem}

\begin{proof}
The analytic dimension inequality gives
\(s_K(x)\le\dim_xX\le2\), so the preceding theorem applies
at every point.  If \(\omega\in\mathscr K_2(U)\), choose
an ordinary cover \(U_i\) and finite local presentations
\(\omega|_{U_i}=\kappa_2(P_i)\).  On each overlap the
germs of \([\alpha_{P_i}]\) and \([\alpha_{P_j}]\) agree
by IQ8.  Since \(\mathscr Q^{1,1}\) is a sheaf, these
classes glue uniquely to a section on \(U\).  Different
local covers and presentations give the same section by the
same germ calculation.  This defines \(\tau_2\).

All formulas are local and real linear, so they commute with
ordinary restriction and addition.  The identity in IQ9 is
checked on the local finite presentations by IQ4.  The proof
glues quotient classes; it does not choose one globally compatible
family of the primitives in IQ8 or one global positive control.
\end{proof}

\begin{corollary}[The fixed-background quadratic ordinary action factors]
\label{cor:r11-intrinsic-quadratic-action-factor}
In the surface sheaf scope above, let \(S\) be an allowed
original background in a legal degree.  The local formula
\[
       \omega=\kappa_2(P)\longmapsto
       \mathfrak A_{S,2}(\omega):=D\bigl(M_S(\alpha_P)\bigr)
                                      =A_S(P)
\tag{IQ10}
\]
defines a real linear ordinary-current action independent of the
local finite presentation.  It is compatible with ordinary
restriction of \(S\) and \(\omega\), and uses every original
compact smooth complementary test and every actual refinement.
The corresponding assertion holds at an arbitrary-dimensional
stalk with \(s_K(x)\le2\).
\end{corollary}

\begin{proof}
For a difference of presentations, IQ8 and the retained QP
first-order scalar-module identity give
\[
 D M_S(\alpha_P-\alpha_{P'})
   =D\{d'(\beta\wedge S)+d''(\zeta\wedge S)\}=0.
\tag{IQ11}
\]
Separate closedness of \(S\), nilpotence and anticommutation
are exactly the hypotheses of that identity.  Thus the local
current formulas agree on overlaps.  Their ordinary current
gluing, or the explicit compact-test partition in SQ, proves
the assertion.  The all-test and refinement quantifiers are
those already proved in OR/QP; no new test class is added.
\end{proof}

\subsection{The finite-union condition cannot be omitted}

At a general higher-residue point, two presentations can each
have finite residue rank at most two while their combined tuple
has larger residue rank.  For example, at the original Gauss
point of \(\mathbb G_{m,K}^3\), the unit tuples
\((u_1,u_2)\) and \((u_2,u_3)\) have residue fields
\(k(\bar u_1,\bar u_2)\) and \(k(\bar u_2,\bar u_3)\),
each of transcendence degree two.  Their union has the field
\(k(\bar u_1,\bar u_2,\bar u_3)\), of transcendence degree
three.  This is a rank diagnostic, not a counterexample to
intrinsic descent.

Accordingly, at an arbitrary original point the comparison of
two specific presentations in IQ8 is proved whenever their
\emph{joint} finite tuple satisfies \(s_f(x)\le2\).  Separate
rank bounds on the two presentations do not justify that step.
The stalk hypothesis \(s_K(x)\le2\) is what makes every finite
union safe and hence gives the canonical map on the entire
quadratic image stalk.  No ordinary open locus is inferred merely
from a pointwise residue-rank bound.

\paragraph{Exact interface boundaries.}
The sheaf morphism IQ9 is a quadratic transgression class modulo
full first-order exactness.  It is not a section of the ordinary
Aeppli cohomology sheaf for a nonzero curvature: its derivative
is the input \(\omega\), not zero.  It is real linear, not a
module map over variable smooth or PS scalar coefficients.  Its
canonical status refers to the displayed finite transgression
formula; no uniqueness among all abstract right inverses of
\(\overline D\) is asserted.

This result proves ordinary restriction and quotient-sheaf locality
for the indicated quadratic image.  It does not construct the full
graded intrinsic differential algebra, impose new functoriality
under arbitrary analytic maps, assert equality on all GK delta
tests, solve arbitrary higher coupled relations or variable-PS/CTO,
or replace any absolute cap carrier.  The local fixed controls
and singular terms remain exactly those supplied by OR/PR.

\paragraph{Source-bound dependencies.}
OR: \texttt{ordinary\_quadratic\_residue\_rank\_two.tex}, SHA-256
\texttt{0a11b646bf6a95afa8b710d2271c0c933fea626fe014e9fde2f1fe7db044f599};
QP: the immutable Round 10
\texttt{fixed\_finite\_ps\_quadratic\_extension.tex}, SHA-256
\texttt{9177fb95c0a2467e06b30ea2829a04c7f165e32abcc9c5d9ac082ec436b9115d}.
The exact existing-sheaf input is the immutable selected QA source
\texttt{qa\_degree\_one\_and\_quadratic\_transgression.tex},
SHA-256
\texttt{a6e8081c94d77b31a4bcf2c5eabd359a105cb03fd5a43f9ec8fe18d2cefdea2d},
especially its opening scope, QA8--QA13, and the proof following
QA11, which already states factorization through the full quotient
\(B^{1,1}/(d'B^{0,1}+d''B^{1,0})\).  The retained R6 full
source at SHA-256
\texttt{fb87839085e500b446d9cc512e5bfca6b8f31de0b8ddf8721d9192d60f9942b0}
identifies this original analytic sheaf as \(B_{\rm Mih}\) in
\texttt{sec:r6-delta-faithfulness}, with its original field and
ordinary topology.  No arbitrary-current faithfulness is used here.

Beyond those existing definitions, the new deduction uses
QP4--QP9, their ordinary restriction statement, OR's full
first-order exactness for the difference relation, and the
ordinary sheaf quotient.  The original analytic proof and
source-bound independent review of OR remain indispensable.
